\documentclass[11pt,a4paper]{article}

\newcommand{\docshorttitle}{HomeSense --- Demonstration Script}
% Shared preamble for the HomeSense project documents.
% Compiled with Tectonic (XeTeX). See build.ps1.

\usepackage[a4paper,margin=25mm,headheight=15pt]{geometry}
\usepackage{fontspec}
\usepackage{amsmath}
\usepackage{amssymb}   % provides \square, used for checklist items
\usepackage{microtype}
\usepackage{booktabs}
\usepackage{tabularx}
\usepackage{longtable}
\usepackage{array}
\usepackage{xcolor}
\usepackage{graphicx}
\usepackage{enumitem}
\usepackage{fancyhdr}
\usepackage{titlesec}
\usepackage{caption}
\usepackage{listings}
\usepackage{tikz}
\usetikzlibrary{positioning,arrows.meta,calc,shapes.geometric,fit,backgrounds}
\usepackage[hidelinks]{hyperref}

% --- typography -------------------------------------------------------------
\setmainfont{Times New Roman}[
  BoldFont={Times New Roman Bold},
  ItalicFont={Times New Roman Italic},
  BoldItalicFont={Times New Roman Bold Italic}
]
\setmonofont{Consolas}[Scale=0.86]

\linespread{1.08}
\setlength{\parskip}{0.6em}
\setlength{\parindent}{0pt}
\setlist{itemsep=0.15em,topsep=0.35em,parsep=0pt}

% --- colours ----------------------------------------------------------------
\definecolor{rule}{HTML}{9AA5B1}
\definecolor{codebg}{HTML}{F5F7F9}
\definecolor{codekw}{HTML}{0B5FA5}
\definecolor{codestr}{HTML}{1B7F3B}
\definecolor{codecom}{HTML}{6B7280}
\definecolor{accent}{HTML}{00696D}
\definecolor{boxbg}{HTML}{F0F4F6}

% --- headings ---------------------------------------------------------------
\titleformat{\section}{\large\bfseries}{\thesection}{0.7em}{}
\titleformat{\subsection}{\normalsize\bfseries}{\thesubsection}{0.7em}{}
\titleformat{\subsubsection}{\normalsize\itshape}{\thesubsubsection}{0.7em}{}
\titlespacing*{\section}{0pt}{1.8em}{0.6em}
\titlespacing*{\subsection}{0pt}{1.2em}{0.4em}

% --- headers and footers ----------------------------------------------------
\pagestyle{fancy}
\fancyhf{}
\renewcommand{\headrulewidth}{0.4pt}
\fancyhead[L]{\small\textsc{\docshorttitle}}
\fancyhead[R]{\small SCS 3311}
\fancyfoot[C]{\small\thepage}

% --- tables -----------------------------------------------------------------
\renewcommand{\arraystretch}{1.25}
\newcolumntype{L}[1]{>{\raggedright\arraybackslash}p{#1}}
\captionsetup{font=small,labelfont=bf,skip=6pt}

% --- inline code ------------------------------------------------------------
% Handles underscores without manual escaping, so field names such as
% max_on_duration can be written naturally.
\newcommand{\code}[1]{\texttt{\small #1}}

% --- code listings ----------------------------------------------------------
\lstdefinestyle{project}{
  backgroundcolor=\color{codebg},
  basicstyle=\ttfamily\scriptsize,
  keywordstyle=\color{codekw}\bfseries,
  stringstyle=\color{codestr},
  commentstyle=\color{codecom}\itshape,
  numberstyle=\tiny\color{codecom},
  breaklines=true,
  breakatwhitespace=true,
  captionpos=b,
  frame=single,
  rulecolor=\color{rule},
  framesep=5pt,
  xleftmargin=8pt,
  xrightmargin=4pt,
  showstringspaces=false,
  tabsize=2,
  columns=fullflexible,
  keepspaces=true,
  upquote=true,
}
\lstset{style=project}

\lstdefinelanguage{Kotlin}{
  keywords={package,as,typealias,class,this,super,val,var,fun,for,while,do,null,
    true,false,is,in,throw,return,break,continue,object,if,else,when,try,catch,
    finally,for,do,while,interface,data,sealed,override,private,public,internal,
    suspend,operator,companion,by,import,enum,const},
  sensitive=true,
  comment=[l]{//},
  morecomment=[s]{/*}{*/},
  morestring=[b]",
}

% --- callout box ------------------------------------------------------------
\newcommand{\keypoint}[1]{%
  \par\vspace{0.4em}%
  \noindent\begin{tikzpicture}
    \node[draw=rule,fill=boxbg,line width=0.4pt,rounded corners=2pt,
          inner sep=9pt,text width=\dimexpr\textwidth-24pt\relax,align=justify]
      {#1};
  \end{tikzpicture}%
  \par\vspace{0.4em}%
}

% --- title page -------------------------------------------------------------
\newcommand{\titlepagestandard}[2]{%
\begin{titlepage}
\centering
\vspace*{2.2cm}

{\Large\scshape University of Colombo School of Computing\par}
\vspace{0.4cm}
{\large SCS 3311 --- Mobile Application Design and Development\par}
\vspace{0.3cm}
{\large Mini-Project\par}

\vspace{2.4cm}
{\rule{\textwidth}{0.8pt}}
\vspace{0.7cm}

{\LARGE\bfseries HomeSense\par}
\vspace{0.5cm}
{\Large Smart Home Monitoring and Control System\par}
\vspace{0.7cm}
{\large #1\par}

\vspace{0.7cm}
{\rule{\textwidth}{0.8pt}}

\vspace{2.4cm}

{\large\bfseries Group Members\par}
\vspace{0.6cm}

\begin{tabular}{@{}L{4.6cm}L{3.1cm}L{6.6cm}@{}}
\toprule
\textbf{Name} & \textbf{Index number} & \textbf{Module} \\
\midrule
R.L. Weerasinghe & 23002204 & Synchronisation, simulator and safety worker \\
W.T. Mahagamage  & 23001038 & Device profiles, control interface and reporting \\
D.M. Isakya      & 23000643 & Floor representation and grid mapping \\
\bottomrule
\end{tabular}

\vfill
{#2\par}
\vspace{0.8cm}
\end{titlepage}
}


% Timing rows are used throughout, so the layout is defined once.
\newenvironment{runsheet}
  {\begin{longtable}{@{}L{1.5cm}L{5.4cm}L{7.9cm}@{}}
   \toprule \textbf{Time} & \textbf{Action} & \textbf{Point to make} \\ \midrule
   \endfirsthead
   \toprule \textbf{Time} & \textbf{Action} & \textbf{Point to make} \\ \midrule
   \endhead}
  {\bottomrule\end{longtable}}

\begin{document}

\titlepagestandard{Demonstration Script}{\large August 2026}

\section*{Purpose and constraints}
\addcontentsline{toc}{section}{Purpose and constraints}

The recording must not exceed \textbf{25 minutes}. This script targets
\textbf{22 minutes}; the margin is deliberate, since live demonstrations
generally run over.

The specification requires all three members to appear, each introducing
themselves and stating their contribution. Each block therefore opens with that
introduction rather than covering all three at the start.

\begin{table}[h]
\centering
\begin{tabular}{@{}L{1.2cm}L{4.6cm}L{8.6cm}@{}}
\toprule
\textbf{Block} & \textbf{Member} & \textbf{Topic} \\
\midrule
1 & D.M. Isakya (23000643)      & Floor representation and device placement \\
2 & W.T. Mahagamage (23001038)  & Device profiles, control interface and reporting \\
3 & R.L. Weerasinghe (23002204) & Synchronisation, simulator and safety worker \\
\bottomrule
\end{tabular}
\end{table}

\section*{Preparation checklist}
\addcontentsline{toc}{section}{Preparation checklist}

\begin{itemize}[label=$\square$]
  \item Worker running (\code{npm start}), with the terminal visible in the frame
  \item Simulator open in a browser, worker status indicator showing online
  \item \code{npm run seed} already executed, so the reporting screen has history
  \item Phone screen mirrored, notifications enabled, do-not-disturb off
  \item Iron's \code{max\_on\_duration} set to 30 seconds
  \item A browser tab with \code{database.rules.json} open for section 6
  \item Screen recorder capturing both the phone and the desktop in one frame
\end{itemize}

\newpage
\section{Introduction --- D.M. Isakya \hfill \normalfont 0:00 -- 1:30}

Introduce yourself with your index number and state your module: floor
representation, the grid, and device placement.

State the problem briefly: a house has several floors, appliances of different
kinds, and some of them present a fire risk if left running.

Then state the central design decision once, clearly:

\keypoint{Every switch in the system stores three separate values: the state the
application requests, the state the hardware reports, and the status the server
derives from the two. That separation is why the error and disconnected states
in this system represent actual observations. The difference will be visible
later in the demonstration.}

\section{Floors and grid --- D.M. Isakya \hfill \normalfont 1:30 -- 5:30}

\begin{runsheet}
1:30 & Floors list showing two storeys with live device counts & Multiple floor plans, each with its own grid \\
2:15 & Open the ground floor & Plans are drawn as geometry rather than images, so they stay sharp at any density and the aspect ratio is exact \\
2:45 & Indicate the legend & Shape encodes device kind, fill pattern encodes status. Nothing relies on colour alone \\
3:15 & Tap an empty cell and place an outlet & Positions are stored as grid cells, not pixels, so a plan renders identically on a phone and a tablet \\
4:00 & Rotate the phone & Placement is unaffected. The coordinate mapping is one pure class with sixteen unit tests \\
4:30 & Add a floor, selecting a plan and grid size & Floors are added and managed at runtime \\
5:00 & Open the gang box & One device, one cell, one heartbeat, with three independently addressable switches. The specification requires a single entity and the data model permits nothing else \\
\end{runsheet}

Hand over to W.T. Mahagamage.

\section{Device profiles --- W.T. Mahagamage \hfill \normalfont 5:30 -- 10:00}

Introduce yourself with your index number: device profiles, the control
interface and usage reporting.

\begin{runsheet}
6:00 & Outlet card, toggle it & The simulator lamp responds within a second, with no refresh anywhere \\
6:45 & Gang box, toggle switch 2 only & Independently addressable; switches 1 and 3 are unaffected \\
7:15 & Master toggle & Addresses each slot in sequence, still as one device \\
7:45 & Schedule editor on the porch light & On at 18:30, off at 06:00 --- a window crossing midnight, which is the usual case \\
8:30 & Safety editor on the iron & \code{max\_on\_duration}, using the specification's own field name. Enforced by the server, not by this screen \\
9:00 & Switch the iron on; the countdown appears & The indicator drains as the permitted time is consumed and changes colour below twenty per cent \\
9:30 & Camera tile, then full screen & Simulated snapshot and stream, as the specification permits. The configured URIs are shown beneath \\
\end{runsheet}

Switch the iron off before continuing; the cut-off is demonstrated separately at
14:00. Hand over to R.L. Weerasinghe.

\section{Simulator and fault injection --- R.L. Weerasinghe \hfill \normalfont 10:00 -- 14:00}

Introduce yourself with your index number: synchronisation, the hardware
simulator and the safety worker. Keep the phone and browser both in frame
throughout this section.

\begin{runsheet}
10:30 & Show the simulator & This represents the hardware. It writes the reported state and a heartbeat, and nothing else. It does not write status \\
11:00 & Physical toggle on the ceiling light & A change originating at the wall switch. The phone was not touched, and the change appears immediately, recorded with source \code{SIMULATOR} \\
11:45 & Simulate fault on the fan, then toggle it from the app & The relay stops responding. Nothing writes an error state \\
12:15 & Wait approximately ten seconds and observe the app & The client reverts its optimistic switch at four seconds and warns; the worker publishes \code{ERROR} at ten seconds. The error was derived by a third process from a disagreement between two others \\
13:00 & Unplug the kitchen outlet & The heartbeat stops. Within fifteen seconds the worker publishes \code{DISCONNECTED} \\
13:30 & Close the simulator tab entirely & The \code{onDisconnect()} handler fires and all devices become disconnected. The absence was recorded by the server \\
\end{runsheet}

Reopen the tab and clear the fault before the next section.

\section{Safety cut-off --- R.L. Weerasinghe \hfill \normalfont 14:00 -- 17:00}

This is the central demonstration.

\begin{runsheet}
14:00 & Set the iron to \code{max\_on\_duration = 30} and switch it on & The countdown begins and the simulator's heat bar fills \\
14:30 & Force-stop the application, showing the settings screen & The phone is now out of the system entirely \\
14:45 & Indicate the worker terminal and the simulator & Only these two processes are running \\
15:00 & Wait, describing the countdown & No timer is held in memory. The worker recomputes elapsed time from the stored \code{onSince} value at every pass \\
15:30 & The cut-off fires: the lamp goes out, a log line appears, a notification arrives & The appliance was switched off with the application force-stopped \\
16:00 & Reopen the app and show the alert centre & The critical alert was written by the server. A client may only mark it acknowledged \\
16:30 & Stop and restart the worker mid-countdown, repeating briefly & A restart loses nothing, because the state is in the database. The next pass acts immediately. This case is covered by a unit test \\
\end{runsheet}

\section{Enforcement and reporting \hfill \normalfont 17:00 -- 19:00}

\begin{longtable}{@{}L{1.5cm}L{3.6cm}L{9.7cm}@{}}
\toprule
\textbf{Time} & \textbf{Member} & \textbf{Action} \\
\midrule
\endfirsthead
17:00 & R.L. Weerasinghe & Show \code{database.rules.json}. A write grant propagates in Realtime Database and cannot be withdrawn at a lower level, but validation rules do restrict client writes, and the Admin SDK bypasses validation. \code{status} therefore carries a validator that accepts a write only when the value is unchanged \\
17:45 & R.L. Weerasinghe & State the limitation directly: the client and simulator share anonymous authentication, so the rules cannot distinguish them. What no client can do is publish a status value, which is the property the safety argument depends on \\
18:00 & W.T. Mahagamage & Usage screen. Every figure derives from logged transitions rather than from comparing snapshots, so the cut-off that occurred with the phone closed is included \\
18:30 & W.T. Mahagamage & Leaderboard, cut-off count and energy estimate, labelled as an estimate because it assumes rated wattage throughout \\
18:50 & W.T. Mahagamage & CSV export through the share sheet \\
\bottomrule
\end{longtable}

\section{Architecture and limitations \hfill \normalfont 19:00 -- 20:30}

One point each:

\begin{itemize}
  \item \textbf{D.M. Isakya} --- MVVM structure: composable, ViewModel,
        repository, data source. No composable accesses Firebase directly.
  \item \textbf{R.L. Weerasinghe} --- Why Realtime Database rather than
        Firestore: per-child listeners and \code{onDisconnect()}. Why a
        standalone worker rather than Cloud Functions: the billing requirement,
        and the fact that the rules are pure functions that would transfer
        unchanged.
  \item \textbf{W.T. Mahagamage} --- What was not done: a single home,
        simulated cameras, estimated energy figures, and the dependency on the
        worker running, which is why the simulator displays the worker's
        heartbeat.
\end{itemize}

\section{Verification --- W.T. Mahagamage \hfill \normalfont 20:30 -- 21:30}

Run both suites on camera:

\begin{lstlisting}[language=bash]
./gradlew test          # 73 tests
cd worker && npm test   # 40 tests
\end{lstlisting}

Highlight two cases: the cut-off fires at exactly \code{max\_on\_duration}, and
it still fires correctly after a simulated worker restart.

\section{Closing --- D.M. Isakya \hfill \normalfont 21:30 -- 22:00}

One sentence each on further work: custom authentication claims to separate
client from hardware, support for multiple homes, and real camera streams.

\newpage
\section*{Priority sequences}
\addcontentsline{toc}{section}{Priority sequences}

If time is short, these three sequences take precedence over the commentary.

\begin{enumerate}
  \item \textbf{14:00 -- 16:00} --- the cut-off firing with the application
        force-stopped.
  \item \textbf{11:00} --- a change originating at the hardware appearing on the
        phone.
  \item \textbf{12:15} --- an error state derived from a sustained disagreement.
\end{enumerate}

\section*{Contingencies}
\addcontentsline{toc}{section}{Contingencies}

\begin{table}[h]
\centering
\begin{tabular}{@{}L{4.4cm}L{10.2cm}@{}}
\toprule
\textbf{Problem} & \textbf{Response} \\
\midrule
Network unavailable & Use the demo build. It runs the same interface against an in-memory backend with a working cut-off, and the fault-injection controls are in the device sheet \\
Worker will not start & Check \code{worker/.env} and the service account key. The simulator's status indicator reports the condition \\
Notification does not arrive & Say so and continue. The cut-off has already occurred and the alert centre confirms it, which is itself the point being made \\
No data appears in the app & Confirm that \code{homeId} matches across the application, the worker's environment file and the simulator configuration \\
\bottomrule
\end{tabular}
\end{table}

\end{document}
